\documentclass[11pt,a4paper]{article}
\usepackage{jheppub}
\usepackage{amsmath,amssymb,booktabs,slashed}
\newcommand{\PL}{P_L}
\newcommand{\dd}{\mathrm{d}}
\preprint{nutau-prop / v0.86 / 31 October 2026}
\title{A decay--scattering normalization audit of resolved neutrino transition operators}
\author{Subhendu Rakshit}
\hypersetup{pdftitle={Resolved neutrino transition normalization audit v0.86},pdfauthor={Subhendu Rakshit}}
\abstract{We resolve a normalization mismatch in the project's implementations of
two PKR-like resolved neutrino transition vertices. For a specified massless,
distinct-state chiral transition, exact traces establish a factor four between
tensor and derivative scattering rates, matching their decay-width ratio.
The inherited derivative and tensor DIS kernels were respectively too large
by factors sixteen and eight relative to the adopted widths. Fresh CT18 folds
and conditional collider screens replace the affected numerical claims.
This is a technical correction to the project implementation, not a finding
that the published PKR calculation is incorrect or a general exclusion theorem.}
\begin{document}
\maketitle

\section{Scope and coefficient conventions}
The resolved interactions considered in ref.~\cite{PKR} motivate the following
explicit vertices for one ordered, distinct-state chiral transition:
\begin{equation}
 V_{44}^\mu=a(k+k')^\mu\PL,\qquad
 V_{43}^\mu=2a q_\nu\sigma^{\nu\mu}\PL,\qquad \kappa=g_q a.
\end{equation}
Here $q=k-k'$, $\PL=(1-\gamma_5)/2$ and
$\sigma^{\mu\nu}=i[\gamma^\mu,\gamma^\nu]/2$, with metric $(+---)$.
The tensor's factor two follows from the field strength. We include the
Hermitian conjugate, with two distinct conjugate invisible decay channels.
The incoming active neutrino has one helicity; only the incoming quark spin
is averaged. Singlet-current colour factors cancel.

The project realization is an active-to-sterile transition. Identifying it
with a literal charge-conjugated active-flavour operator requires additional
field, flavour and multiplicity information. Neither a sterile interpretation
nor the original UFO conventions follow from the equation number alone.

An additional exact selection rule matters. For anticommuting fields,
$(C\PL)^T=-C\PL$ and
$(C\sigma^{\mu\nu}\PL)^T=+C\sigma^{\mu\nu}\PL$ imply that both the
charge-conjugated tensor bilinear and the scalar bilinear with an
$\overleftrightarrow{\partial}$ derivative are antisymmetric in flavour.
They vanish for literally identical diagonal fields. A nonzero derivative
realization also needs distinct fields or an off-diagonal assignment;
suppressed printed indices do not specify the original numerical model.

\section{Analytic scattering and crossing checks}
For $\nu(k)q(p)\to N(k')q(p')$, write $s=\hat s$, $t=(k-k')^2$,
$u=-s-t$ and $K=k+k'$. Up to a phase,
\begin{equation}
 \mathcal M_{44}=\frac{\kappa}{t-M^2}
 [\bar u(k')\PL u(k)]K_\mu[\bar u(p')\gamma^\mu u(p)].
\end{equation}
Current conservation and $q\cdot K=0$ eliminate the longitudinal propagator.
The traces give
\begin{align}
 S&=\operatorname{Tr}[\slashed k'\PL\slashed k P_R]=-t,\\
 H^{\mu\nu}&=\tfrac12\operatorname{Tr}[\slashed p'\gamma^\mu\slashed p\gamma^\nu],
 &K_\mu K_\nu H^{\mu\nu}&=-4su,\\
 \overline{|\mathcal M_{44}|^2}&=\frac{4\kappa^2stu}{(M^2-t)^2}.
\end{align}
Using $t=-sy$ and $u=-s(1-y)$,
\begin{equation}
 \frac{\dd\hat\sigma_{44}}{\dd y}
 =\frac{\kappa^2s^2y(1-y)}{4\pi(M^2+sy)^2}.
\end{equation}
The massless external-fermion equations relate the tensor current to twice
the derivative current up to a phase. This relation holds for virtual
$t$-channel exchange: it does not assume an on-shell mediator. Therefore
\begin{equation}
 \dd\sigma_{43}=4\dd\sigma_{44}
 \quad\text{at equal }a,g_q,M\text{ and phase-space domain}.
\end{equation}
For $s_N=2m_NE$, $Q^2=xs_Ny$ and $F=\sum_q(xf_q+xf_{\bar q})$,
\begin{equation}
 \frac{\dd^2\sigma_{44}}{\dd x\dd y}
 =\frac{\kappa^2x s_N^2 y(1-y)F(x,Q^2)}{4\pi(M^2+Q^2)^2}.
 \label{eq:fold}
\end{equation}
The contact limit has $\hat\sigma_{44}=\kappa^2s^2/(24\pi M^4)$,
with relative corrections of order $s/M^2$.

The vector decay polarization average is $1/3$. The one-channel squared
amplitudes are $a^2M^4/3$ and $4a^2M^4/3$. Including both conjugate channels,
\begin{equation}
 \Gamma_{44}=\frac{a^2M^3}{24\pi},\qquad
 \Gamma_{43}=\frac{a^2M^3}{6\pi}.
\end{equation}
These agree with the project's physical widths. In particular,
\begin{equation}
 \frac{1}{\Gamma_{\rm inv}}\frac{\dd^2\sigma}{\dd x\dd y}
 =\frac{6g_q^2}{M^3}\frac{x s_N^2y(1-y)F}{(M^2+Q^2)^2}
 \label{eq:invariant}
\end{equation}
is identical for both vertices at this multiplicity. A coefficient relabelling
cannot repair a mismatch in this invariant. Extra invisible species require
an explicit change in the width-to-incident-species relation.

The inherited derivative DIS coefficient was $4/\pi$ instead of $1/(4\pi)$.
The inherited tensor kernel was twice that value, while its width was four
times the derivative width. Relative to the specified widths these kernels
overstated the rates by sixteen and eight. The previous width calculation
survives; the previous cross-section identification does not.

\section{Numerical checks and corrected benchmarks}
We used CT18NLO member zero with a leading-order hard process, an isoscalar
target, five quark flavours, $Q\geq1.3$ GeV and $x\geq10^{-9}$.
The integration starts at the physical $Q_{\min}$-dependent $y$ boundary.
The reported ratio is the mean of $\sigma_{\rm NP}/\sigma^{\nu}_{\rm SM}$
and $\sigma_{\rm NP}/\sigma^{\bar\nu}_{\rm SM}$, with SM including CC and NC.
It is not a ratio to NC alone or a detector event ratio.

\begin{table}[ht]
\centering
\begin{tabular}{lrrrr}
\toprule
Benchmark & $g_q$ & $a\ (\mathrm{GeV}^{-1})$ & $R(100\,\mathrm{PeV})$ & $R(220\,\mathrm{PeV})$\\
\midrule
v0.73, Eq. 4.4 & .030 & .00356307 & 0.3106\% & 0.3046\%\\
Early Eq. 4.3 & .030 & .00197258 & 0.6217\% & 0.6204\%\\
v0.81, Eq. 4.4 & .038 & .00237978 & 0.3179\% & 0.3155\%\\
v0.83, Eq. 4.3 & .017 & .00674371 & 0.6212\% & 0.5972\%\\
\bottomrule
\end{tabular}
\caption{Fresh fixed-coupling results at $M=50$ GeV and $\sqrt{\hat s}<1/a$.
The first two 100-PeV NP cross sections are $2.17770\times10^{-35}$ and
$4.35948\times10^{-35}\ \mathrm{cm}^2$.}
\end{table}

Changing the cutoff to $0.5/a$ at fixed couplings lowers the first two
100-PeV ratios to $0.1232\%$ and $0.3416\%$. Since $a=c/\Lambda$, the choice
$1/a$ is a convention for the Wilson coefficient, not a universally established
UV threshold. Equal amplitudes require $a_{44}=2a_{43}$. At a common physical
cutoff both rates and widths then agree; imposing $1/a$ separately gives
different domains and a prescription-dependent hierarchy. This is not a proof
of equivalence between complete electroweak or UV theories.

Exact symbolic matrix traces pass at general scattering angle and arbitrary
decay mass, including the Ward identities. Twenty tests check the contact
limit, endpoints, coefficient mapping and analytic coupling bounds. Fourteen
Mac Studio worker processes, each with single-thread numerical libraries,
computed 288 base folds and targeted refinements. The largest relative
quadrature change at refined maxima was $1.42\times10^{-7}$, and the largest
interpolation error was $1.93\times10^{-4}$. These figures do not include
PDF, perturbative, low-$Q$ or experimental uncertainties.

\section{Conditional collider screen}
At fixed mass and in the narrow-width approximation, let $z=g_q^2$,
$I=\Gamma_{\rm inv}$, $B=\Gamma_{qq}/g_q^2$ and
\begin{equation}
 P=\frac{zI}{Bz+I},\qquad r=K_M P,\qquad D=z-P=g_q^2\mathrm{BR}_{qq}.
\end{equation}
The bounds on $z$ can be solved analytically. A monojet condition
$P\leq P_{\max}$ gives
$z\leq P_{\max}I/(I-BP_{\max})$ only if $I>BP_{\max}$.
A dijet condition $D\leq d$ gives
\begin{equation}
 z\leq\tfrac12\left[d+\sqrt{d^2+4dI/B}\right].
\end{equation}
We retain $\Gamma/M\leq0.1$, $g_q^2\leq4\pi$ and $1/a>M$.
Every sampled $a$ interval was optimized and its winning value checked by a
direct fold, avoiding a first-crossing-only scan.

CMS reports universal-vector dijet coupling limits from 0.03 to 0.13 over
50--300 GeV, assuming a fully hadronic decay~\cite{CMS}. We deliberately use
the weakest value, $g_q\sqrt{\mathrm{BR}_{qq}}\leq0.13$, at each tested mass.
This relaxes the mass-dependent constraint under the narrow-width mapping.
The resulting nominal $r\leq1$ maxima are:
\begin{center}
\begin{tabular}{rrr}
\toprule
$M\ (\mathrm{GeV})$ & Eq. 4.4 & Eq. 4.3\\
\midrule
50 & 3.142\% & 4.514\%\\
100 & 0.2197\% & 0.2635\%\\
220 & 0.0583\% & 0.0743\%\\
\bottomrule
\end{tabular}
\end{center}
These are conditional screens, not exact 95\% CL exclusions. The $M=50$ GeV
normalization uses the existing 100k-event same-seed reference; 100 and 220 GeV
use 20k-event pilot shapes. Exact mass-specific dijet data were not recovered.
With the same deliberately weak dijet envelope but $r\leq1.2$, Eq. 4.3 at
50 GeV reaches 5.258\%. Thus the relaxed nominal bound must not be promoted
to a robust exclusion. Approximate mass-specific inputs yield smaller effects,
but remain sensitivity diagnostics. Other mediator masses are not closed.

At the same couplings, correcting DIS leaves the existing collider $r$
unchanged because the widths already had this normalization. What changes is
whether that point satisfies the UHE target. The $r=1.239$ benchmark is no
longer a 5\% UHE benchmark.

\section{Propagation and interpretation}
For the specified sterile final state there is no SM interference and the
new transport term is an active-flux sink. The unregenerated optical depth
$N_AX\sigma_{\rm NP}$ inherits the rate correction. An already regenerated
SM spectrum cannot simply be multiplied by a sink exponential: NC downscattering
and tau decay couple energies. We do not claim a new transport solution or a
tau-specific IceCube likelihood here.

The corrected benchmark effects remain subpercent at 220 PeV. A large
relative change on an opaque vertical chord is not, by itself, an observable
cosmogenic event enhancement. The original objective remains a valid model
with measurable tau-neutrino propagation distortions after laboratory,
IceCube, validity and detector constraints. These benchmarks do not establish
that goal, but neither do three mass screens prove a general no-go theorem.

The error established here belongs to the project's normalization chain.
An attribution to the published PKR analysis requires reproducing its actual
field conventions and collider model. Modern experimental constraints and
implementation errors must be discussed separately. All superseded results
are preserved, with corrected code and machine-readable outputs in v0.86.

\paragraph{The concurrently added published-bound reconstruction.}
PKR quotes 5-GeV IceCube bounds $\kappa_{43}=4.8\times10^{-4}$ and
$\kappa_{44}=2.5\times10^{-4}\ \mathrm{GeV}^{-1}$~\cite{PKR}. If the two
bounds use an identical allowance, flavour sum and the coefficient convention
above, they imply $F_{43}/F_{44}=(2.5/4.8)^2=0.271267$ for
$\sigma_i=\kappa_i^2F_i$, instead of four. Our kernels evaluated at those
couplings have ratio 14.7456. This is a conditional normalization discrepancy,
but the pair does not determine which absolute implementation is wrong.
Trusting the 4.3 anchor gives a 4.4 bound of $9.6\times10^{-4}$; trusting the 4.4
anchor gives a 4.3 bound of $1.25\times10^{-4}\ \mathrm{GeV}^{-1}$.
Neither is a new IceCube fit. The concurrent v0.85 note's inferred coefficient
58.98 has rounding range 55.5--62.7, which does not contain 64. Extra uncertainty
is necessary to support that suggested identification. This distinction
separates the directly demonstrated repository error from a hypothesis about
the published implementation.

\begin{thebibliography}{9}
\bibitem{PKR} S.~Pandey, S.~Karmakar and S.~Rakshit,
\emph{Strong constraints on non-standard neutrino interactions: LHC vs. IceCube},
\href{https://arxiv.org/abs/1907.07700}{arXiv:1907.07700}.
\bibitem{CMS} CMS Collaboration,
\emph{Search for low-mass vector and scalar resonances decaying into a
quark-antiquark pair in proton-proton collisions at $\sqrt{s}=13$ TeV},
CMS-EXO-24-007, \href{https://arxiv.org/abs/2603.21965}{arXiv:2603.21965};
\href{https://cms-results.web.cern.ch/cms-results/public-results/publications/EXO-24-007/}{official results and branching-ratio assumptions}.
\end{thebibliography}
\end{document}
